\documentclass[8pt]{beamer}

\setbeamertemplate{bibliography item}{\insertbiblabel}

\setbeamersize{text margin left=5mm,text margin right=5mm}

\usepackage{/Users/wdlinch3/Dropbox/Rashoumon/LaTeX/macros/WDL3beamer}

%Pausing toggle
\def\mypause{
%\pause
} %Remove "%" or comment out 
% Pro-tip: for 50+10 minute talk there should < about 50 pauses

%%%%%%%%%%%%%%%%%%%%%%%%%%%%%%%%%%%%%%%%%%%%%%%%%%%%%%%%%%%%%%%%
%%%%%%%%%%%%%%%%%%%%%%%%%%%%%%%%%%%%%%%%%%%%%%%%%%%%%%%%%%%%%%%%
%%%%%%%%%%%%%%%%%%%%%%%%%%%%%%%%%%%%%%%%%%%%%%%%%%%%%%%%%%%%%%%%
%%%%%%%%%%%%%%%%%%%%%%%%%%%%%%%%%%%%%%%%%%%%%%%%%%%%%%%%%%%%%%%%
\usepackage{beamerthemesplit}

\title{Supergeometry of Gauged $p$-Form Hierarchies}
%\author{William D. Linch III \\
%~\\
%{\em George P. and Cynthia W. Mitchell Institute for\\
%Fundamental Physics and Astronomy, \\
%Texas A\&{}M University.}
%}
\subtitle{(Based on work with M. Becker, K. Becker, and D. Robbins)}
\author{William D. Linch III}
\institute{ 
\includegraphics[height=5mm]{TAMUlogo.jpg}\\
\em Mitchell Institute for Fundamental Physics and Astronomy, \\
Texas A\&{}M University
}

\date{
Prepared for the FRG %Focused research group
Workshop on \\
{\em 
Special Holonomy Geometry, Mirror, and Supersymmetry, 
\\
Harvard University\\
%~\\
9|10 May, 2016}
}


%\logo{\includegraphics[height=5mm]{TAMUlogo.jpg}}

\begin{document}

\frame{\titlepage}

%\section[Outline]{}
%\frame{\tableofcontents}

%%%%%%%%%%%%%%%%%%%%%%%%%%%%%%%%%%%%%%%%%%%%%%%%%%%%%%%%%%%%%%%%
%%%%%%%%%%%%%%%%%%%%%%%%%%%%%%%%%%%%%%%%%%%%%%%%%%%%%%%%%%%%%%%%
\section{Introduction}
%%%%%%%%%%%%%%%%%%%%%%%%%%%%%%%%%%%%%%%%%%%%%%%%%%%%%%%%%%%%%%%%
\frame{
\frametitle{Motivation: Eleven-dimensional Supergravity}
Investigations into M(ystery)-theory: Fundamental (M-branes) and/or \underline{Effective (11D SG)}\\~\\
%
\mypause
%
11D SG supermultiplet: $g$ (metric), $\psi\in \Omega^1(S)$ (gravitino), $C\in \Omega^3$\\
Action ($G=dC$): 
\begin{align*}
S \sim 
	S_{EH}
	+ \tfrac i2 \int d^{11}x \,  \bar \psi_m \Gamma^{mnp} \nabla_n \psi_p 
	+ \tfrac12 \int G\wedge \ast \, G
	+ \int C \wedge G\wedge G
\end{align*}
Looks innocuous but ...
%
\mypause
%
\begin{enumerate}
\item Supersymmetry mixes metric and 3-form components
\item Supersymmetry algebra does not close on gravitino unless we impose $\Gamma^{mnp} \nabla_n \psi_p =0$
\begin{center}
	mixed {\em kinematical} (first class/off-shell) and {\em dynamical} (second class/on-shell)\\
	$\Rightarrow$ supersymmetry algebra is ``soft'' (closes only on-shell) \\
\end{center}
\end{enumerate}
%
\mypause
%
Super(symmetric)space does not help because we need 
\begin{enumerate}
\item Supersymmetric space as Klein geometry $G/H$ with $G=$ eleven-dimensional super-Poincar\'e group and $H=Spin(10,1)$ Lorentz subgroup.
\item Maximally non-integrable odd distribution $\{D\}$: $[D,D] = -i \!\not\! \partial $ and $\not\! \partial^2=\Box$
\item ``Large enough'' integrable odd sub-distribution $\{\bar D\}$ ($\nexists$ in 11D)
\end{enumerate}

} %end frame

%%%%%%%%%%%%%%%%%%%%%%%%%%%%%%%%%%%%%%%%%%%%%%%%%%%%%%%%%%%%%%%%
\frame{
\frametitle{Goal}
Suggested solution: 
\begin{enumerate}
\item Reduce $H\to H_{st} \times H_{int}$ (and $G\to G_{st} \times G_{int}$) such that Klein geometry $(G_{st},H_{st})$ admits integrable odd sub-distribution
\item Replace 11D superspace by foliation by $G_{st}/H_{st}$ superspaces
\end{enumerate}
%
\mypause
%

{\bf Goal:} \\	
Understand 11D SG in terms of foliation by 4D, $N=1$ ``simple'' superspaces \\
$\Rightarrow$ finite-dimensional off-shell representations encode kinematical part of symmetry\\
Remaining dynamical part of symmetry encoded by additional structures:
%
\mypause
%
%{\bf Goal:} \\	
\begin{itemize}
\item $\mathrm D>4$ and/or $N>1$ supergravity contains graviton and gravitino but also
$p$-form fields + Chern-Simons-like terms ({\it e.g.}\ 5D, $N=1$ $\Rightarrow$ $p=1$ and 11D, $N=1$ $\Rightarrow$ $p=3$)
%
\mypause
%
\item Need to understand $p$-forms in $D>4$ in terms of 4D, $N=1$\\ 
$\Rightarrow$ ``hierarchies'' of 4D superforms
%
\mypause
%
\item Need to understand associated Chern-Simons-like superinvariants
%
\mypause
%
\item Superforms formulated in terms of frames \\
$\Rightarrow$ Reduction to 4D gives Kaluza-Klein gauge fields $e_m{}^i$\\ 
$\Rightarrow$ non-abelian gauging of $p$-form hierarchy
\end{itemize}
} %end frame

%%%%%%%%%%%%%%%%%%%%%%%%%%%%%%%%
\frame{
\frametitle{$N=0$}
\begin{itemize}
\item Spacetime $M$ with $p$-form gauge field $\phi_{[p]}$ transforming under abelian gauge transformation as $\delta \phi_{[p]} = d_M \Lambda_{[p-1]}$, with invariant field strength $F_{[p+1]} = d_M \phi_{[p]}$, and Chern-Simons-like invariant 
\begin{align*}
S_{CS} =  \int_M \phi \wedge d\phi\wedge \dots \wedge d\phi = \int_M \phi \wedge c
\end{align*}
%%
%\mypause
%%
\item Splitting $M = X\times Y$ $\Rightarrow$ $d_M=d_X+d_Y$ and $\Omega^p_M = \bigoplus_{q=0}^p \Omega_X^{p-q}(\Omega_Y^q)$ double complex.
%with $X\sim \mathbf R^4$ and $Y\sim \mathbf R^7$ 
%%
%\mypause
%%
\item Write $d:=d_X$ and $\partial := d_Y$. Forms split $\phi_{[p]} \to \sum_q \phi_{[p-q,q]}$ by bi-degree and 
\begin{align*}
\delta \phi_{[r,s]} = d \Lambda_{[r-1,s]} + \partial \Lambda_{[r,s-1]}
~~~\mathrm{and}~~~
F_{[r,s]} = d \phi_{[r-1,s]} + \partial \phi_{[r,s-1]}
\end{align*}
%%
%\mypause
%%
\item Chern-Simons-like action for $\{\phi_{[r,s]}\}_{r+s=p}$:
\begin{align*}
S_{CS}  = \sum_{r=0}^p%{\mathrm{dim}(X)}%\sum_{s=0}^{\mathrm{dim}(Y)}
	\int_{X\times Y} \phi_{[r,p-r]} \wedge c_{[\mathrm{dim}(X)-r,\mathrm{dim}(Y)-p+r]}
\end{align*}
%%
%\mypause
%%
\item Invariant (up to surface terms) provided $c$ is a closed: $dc_{[r,s]} + (-1)^p\partial c_{[r+1,s-1]}= 0$
\end{itemize}

} %end frame


%%%%%%%%%%%%%%%%%%%%%%%%%%%%%%%%%%%%%%%%%%%%%%%%%%%%%%%%%%%%%%%%
%%%%%%%%%%%%%%%%%%%%%%%%%%%%%%%%%%%%%%%%%%%%%%%%%%%%%%%%%%%%%%%%
\section{Superspace}

%%%%%%%%%%%%%%%%%%%%%%%%%%%%%%%%%%%%%%%%%%%%%%%%%%%%%%%%%%%%%%%%
\frame{
\frametitle{4D, $N=1$ Superspace}
\begin{itemize}
\item ``Physical'' superspace: Klein geometry $X = G/H$ with $G=$ super-Poincar\'e (4D, $N=1$) and $H\subset G$ Lorentz subgroup $Spin(3,1)\cong SL_2(\mathbf C)$
%
\mypause
%
\item $X \cong \mathbf R^{4|4}$ with Cartesian coordinates $\{z^M\} = \{x^m, \theta^\mu, \bar \theta^{\dt\mu}\}$ \\
Super-translations generated by $\{Q_A\}=\{P_a, Q_\alpha, \bar Q_{\dt \alpha}\}$
%
\mypause 
%
\item Left-invariant framing of the tangent space $[Q , D]=0$:
\begin{align*}
D_A  = e_A{}^M \partial/\partial z^M = 
\left\{
\begin{array}{l}
D_a =  \partial_a := \delta_a^m \partial/\partial x^m \\
D_\alpha = \delta_\alpha^\mu \partial/\partial \theta^\mu 
	-i (\sigma^m)_{\alpha \dt \alpha} \delta_{\dt \mu}^{\dt \alpha} \bar\theta^{\dt \mu}\partial/\partial x^m\\
\bar D_{\dt \alpha} = \delta^{\dt \mu}_{\dt \alpha} \partial/\partial \theta^\mu  	
	-i (\sigma^m)_{\alpha \dt \alpha} \delta_{\mu}^{\alpha} \theta^{\mu}\partial/\partial x^m
\end{array}
\right.
\end{align*}
%
\mypause
%
\item 
``$D$-algebra'' rules: %$[D,D] = 0$, $[D,\bar D] = -2i \partial$, and $[\bar D,\bar D] = 0$
\begin{align*}
[D,D] = 0
~~~,~~~
[D,\bar D] = -2i \partial 
~~~,~~~
[\bar D,\bar D] = 0 
\end{align*}
%
\mypause
%
\item $\Rightarrow$ $\{\bar D\}$ defines integrable distribution 
$\Rightarrow$ ring of ``chiral'' superfields $\{\Phi : \bar D\Phi = 0\}$\\
$\Rightarrow$ off-shell, irreducible superfield representations for any superspin
%
\mypause 
%
\item Super-de Rham differential $d = dz^M \partial/\partial z^M \to e^A D_A$ but co-frame carries torsion $de^A = T^A$ (even in flat superspace) $\Rightarrow$ complicates analysis of closed super-$p$-forms
\end{itemize}
} %end frame


%%%%%%%%%%%%%%%%%%%%%%%%%%%%%%%%%%%%%%%%%%%%%%%%%%%%%%%%%%%%%%%%
\frame{
\frametitle{Closed Superforms}
Solving $dF=0$ in superspace (a.k.a.\ solving the ``Bianchi identities'') gives
\begin{align*}
\begin{array}{ccl}
p=0 
	&& F = f
\\
p=1
	&& F = e^\alpha D_\alpha F + e^a (D\sigma_a \bar D) F
\\
p=2
	&& F = e^\alpha e^a (\sigma_a \bar W)_\alpha + e^a e^b (D \sigma_{ab}W) + \mathrm{h.c.}
\\
p=3
	&& F = e^\alpha e^{\dt \alpha} e^a (\sigma_a)_{\alpha \dt \alpha}H 
		+ \dots 
		+e^a e^b e^c \epsilon_{abcd} (D\sigma^d \bar D)H 
\\
p=4
	&&F = e^\alpha e^\beta e^a e^b (\sigma_{ab})_{\alpha \beta} G
		+\dots
		+ e^a e^b e^c e^d\epsilon_{abcd} D^2G
\end{array}
\end{align*}
%
\mypause 
%
together with the conditions on the superfield strengths
{\renewcommand{\arraystretch}{1.5} %adds some padding
\begin{align*}
\begin{array}{c|crcrclcc}
p
	&& \textrm{BI}
	&& \textrm{irreducibility}
	&& \textrm{solution}
	&& \textrm{prepotential}
\\
\hline
0 
	&&\partial f =0 
	&& \bar f = - f
	&\Rightarrow& f= i c
	&& c \in \mathbf R
\\
1 
	&&\bar D^2 D F =0 
	&& \bar F = F
	&\Rightarrow& F=\tfrac1{2i} \left( \Phi - \bar \Phi	\right)
	&& \bar D\Phi = 0
\\
2
	&& DW- \bar D \bar W = 0%\mathrm{Im}(DW)=0
	&&\bar DW =0 
	&\Rightarrow& W= -\tfrac14 \bar D^2D V
	&& \bar V = V
\\
3
	&&\bar D^2 H =0 
	&& \bar H = H
	&\Rightarrow& H= \tfrac1{2i} \left( D\Sigma - \bar D \bar \Sigma 	\right) %\mathrm{Im}(D\Sigma)
	&& \bar D\Sigma = 0
\\
4
	&&
	&&\bar DG =0 
	&\Rightarrow& G=-\tfrac14 \bar D^2 X
	&& \bar X = X
\\
\textrm{``$5$''}
	&&
	&&
	&&\hfill \Gamma\hfill
	&& \bar D\Gamma = 0\\
\end{array}
\end{align*}
} %end arraystretch
%Organize into ``prepotential'' complex
%\begin{align*}
%%\Omega^\bullet : 
%0\longrightarrow \underbrace{{\mathscr O^0}}_{\textrm{chiral scalar}}
%		\stackrel {d_0} \longrightarrow \underbrace{{\mathscr O^1}}_{\textrm{real scalar}}
%		\stackrel {d_1} \longrightarrow \underbrace{{\mathscr O^2}}_{\textrm{chiral spinor}}
%		\stackrel {d_2} \longrightarrow \underbrace{{\mathscr O^3}}_{\textrm{real scalar}}
%		\stackrel {d_3} \longrightarrow \underbrace{{\mathscr O^4}}_{\textrm{chiral scalar}}
%%		\stackrel {d_4} \longrightarrow \mathscr O^5
%\longrightarrow 0
%\end{align*}
} %end frame

%%%%%%%%%%%%%%%%%%%%%%%%%%%%%%%%%%%%%%%%%%%%%%%%%%%%%%%%%%%%%%%%
\frame{
\frametitle{Prepotential de Rham Complex}
Organize into (exact) cochain complex of sheaves of superfields
\begin{align*}
%\Omega^\bullet : 
\hspace{-5mm}
0\longrightarrow 
i\mathbf R\longrightarrow
		\underbrace{{\mathscr O^0}}_{\textrm{chiral scalar}}
		\stackrel {d_0} \longrightarrow \underbrace{{\mathscr O^1}}_{\textrm{real scalar}}
		\stackrel {d_1} \longrightarrow \underbrace{{\mathscr O^2}}_{\textrm{chiral spinor}}
		\stackrel {d_2} \longrightarrow \underbrace{{\mathscr O^3}}_{\textrm{real scalar}}
		\stackrel {d_3} \longrightarrow \underbrace{{\mathscr O^4}}_{\textrm{chiral scalar}}
		\stackrel {d_4} \longrightarrow 0
\end{align*}
%
\begin{align*}
d_0 \phi = \tfrac1{2i} \left( \phi - \bar \phi	\right) %\mathrm{Im}(\phi)
~,~
d_1 \phi = -\tfrac14 \bar D^2D \phi
~,~
d_2 \phi = \tfrac1{2i} \left( \iota_D \phi  - \iota_{\bar D}\bar \phi\right)
~,~
d_3 \phi = -\tfrac14 \bar D^2 \phi
~,~
d_4 \phi = \bar D\phi
\end{align*}
%
\mypause 
%
\begin{align*}
\textrm{$D$-algebra rules} 
~~~\Rightarrow~~~
d_{i+1}\circ d_i = 0 
\end{align*}
%
\mypause 
%
For $\mathbf A = (\Phi,V ,\Sigma, X, \Gamma) \in \mathscr O^\bullet$ the poly-superform $\mathbf F = d \mathbf A =(0,F, W, H , G)$ is invariant under $\delta \mathbf A = d \bm \Lambda$.
%
\mypause 
%
~\\
~\\
{\bf NB}: These are not derivations:
$d_i (\phi \phi') \neq (d_i\phi) \phi' + \phi d_i \phi'$
\\~\\
\centering
{\em There is no differential graded algebra structure at the level of prepotentials}.
} %end frame

%%%%%%%%%%%%%%%%%%%%%%%%%%%%%%%%%%%%%%%%%%%%%%%%%%%%%%%%%%%%%%%%
\frame{
\frametitle{Simple Superspace Foliation}
Take spacetime $M=X\times Y$ with $X=\mathbf R^{4|4}$ and $Y$ an ordinary $N=0$ space.\\
Extend all $\mathscr O_X^\bullet$ to take values in the de Rham complex $(\Omega_Y^\bullet, d_Y)$; denote $d_Y= \partial$

%
\mypause 
%
Example: 11D SG has $Y=\mathbf R^7$ and 3-form decomposing under $\Omega_M^3 \to \bigoplus_{q=0}^3 \Omega_X^{3-q}(\Omega_Y^q)$ :
{\renewcommand{\arraystretch}{1.2} %adds some padding
\begin{align*}
C_{[3]} & \longrightarrow 
	\begin{array}{ccccccc}
	1 && 7 && {7\choose2} = 21&& {7\choose3} = 35\\
	\textrm{3-form} && \textrm{2-forms} && \textrm{vectors}&& \textrm{scalars}\\
	C_{[3,0]} &+& C_{[2,1]} &+& C_{[1,2]} &+& C_{[0,3]} \\
	X_{\bm 1} && \Sigma_{\bm 7} && V_{\bm {21}} && \Phi_{\bm {35}}
	\end{array}\cr
F_{[4]} & \longrightarrow 
	\begin{array}{cccccccccc}
	F_{[4,0]} &+& F_{[3,1]} &+& F_{[2,2]} &+& F_{[1,3]}&+& F_{[0,4]} \\
%	\textrm{4-form} && \textrm{3-forms} && \textrm{2-forms}&& \textrm{}\\
	G_{\bm 1} && H_{\bm 7} && W_{\bm {21}} && F_{\bm {35}} && E_{\bm {35}}
	\end{array}
\end{align*}
}
%
\mypause 
%
Field strengths and Bianchi identities (relative closure relations):
\begin{align*}
\begin{array}{ccc}
\begin{array}{ccccccc}
E &=& &-& \partial \Phi\\
F &=& d_0(\Phi) &-& \partial V \\ 
W&=& d_1(V) &-& \partial \Sigma \\ 
H&=& d_2(\Sigma) &-& \partial X \\ 
G&=& d_3(X) 
\end{array}
&\Longrightarrow&
\begin{array}{ccc}
d_0(E)&=&- \partial F\\ 
d_1(F) &=& - \partial W \\ 
d_2(W) &=& - \partial H \\
d_3(H) &=& - \partial G \\ 
d_4(G) &=& 0
\end{array}
\end{array}
\end{align*}
%
\mypause 
%
Manifestly supersymmetric ``Abelian Tensor Hierarchy''
%(here ``tensor'' means $p$-form)
} %end frame

%%%%%%%%%%%%%%%%%%%%%%%%%%%%%%%%%%%%%%%%%%%%%%%%%%%%%%%%%%%%%%%%
\frame{
\frametitle{Abelian Chern-Simons-like Action}
``Abelian Tensor Hierarchy'' has Chern-Simons-like invariants of the form $(\textrm{prepotential})\times(\textrm{superfield strength})\times\dots \times(\textrm{superfield strength})$
%\\
%%
%\mypause 
%%
%In 11D $\to$ 4D superspace case we have (anti-)chiral terms and real ones
\begin{align*}
S_{CS} = \int d^4 x \int_Y \Bigg\{
	\int d^2 \theta \left[ \Phi \wedge c_{[4,4]} + \Sigma^\alpha \wedge c_{[2,6] \alpha}\right]
	+\mathrm{h.c.}
\cr
+\int d^4 \theta \left[ V\wedge c_{[3,5]} + X\wedge c_{[1,7]}\right]
\Bigg\}
\end{align*}
%
\mypause 
%
Gauge invariant (up to surface terms) under 
\begin{align*}
\delta (\Phi, V, \Sigma, X) = (0,d_0\Lambda, d_1U, d_2 \Upsilon)  +\partial (\Lambda, U , \Upsilon,0 )
\end{align*}
provided the $c$s satisfy the ``descent relations'' (cocycle condition)
\begin{align*}
\begin{array}{ccc}
0&=&\partial c_{[1,7]}\\ 
d_1(c_{[1,7]}) &=& \partial c_{[2,6]} \\ 
d_2(c_{[2,6]}) &=& \partial c_{[3,5]} \\ 
d_3(c_{[3,5]}) &=& \partial c_{[4,4]} \\
d_4(c_{[4,4]}) &=& 0
\end{array}
\end{align*}
Formally the same equations as the Bianchi identities... \\ 
%
\mypause 
%
but far less trivial because $c$'s are bilinear in field strengths and $d_i$ are not derivations
} %end frame


%%%%%%%%%%%%%%%%%%%%%%%%%%%%%%%%%%%%%%%%%%%%%%%%%%%%%%%%%%%%%%%%
\frame{
\frametitle{Existence of Quadratic Co-cycle}
Explicit solution to the quadratic descent relations
%
\mypause 
%
\begin{align*}
\mathscr O^1 (\Omega_Y^7) \ni c_{[1,7]} &= \mathrm{Re}(E) \wedge F
\cr
\mathscr O^2 (\Omega_Y^6) \ni c_{[2,6]} &= - E  \wedge W -\tfrac i4 \bar D^2 (F  \wedge D F)
\cr
\mathscr O^3 (\Omega_Y^5) \ni c_{[3,5]} &= - \mathrm{Re}(E)  \wedge H -2 \omega (W, F)
\cr
\mathscr O^4 (\Omega_Y^4) \ni c_{[4,4]} &= \boxed{E G} +  \boxed{2W \wedge W} -\tfrac i2 \bar D^2 (FH)
\end{align*}
%
\mypause 
%
Here $\omega (W, F)$ is the ``Chern-Simons superform'' 
\begin{align*}
\omega (W, F) := W^\alpha \wedge D_\alpha F + \tfrac12 D^\alpha W_\alpha \wedge  F +\mathrm{h.c.}
\end{align*}
It is the superspace analogue of $A\wedge F$ satisfying 
\begin{align*}
d_3 \omega (W, F) &= W \wedge d_1 F +  d_2(W) \wedge  F  
\cr
&= W \wedge \partial W +  \partial H \wedge  F  
\end{align*}
%
\mypause 
%
Using every $D$-identity, Bianchi identity, and internal de Rham DGA property, it is possible to show that $c \in \bigoplus_i \mathscr O^{8-i}(\Omega_Y^i)$ is a cocycle.
} %end frame


%%%%%%%%%%%%%%%%%%%%%%%%%%%%%%%%%%%%%%%%%%%%%%%%%%%%%%%%%%%%%%%%
%%%%%%%%%%%%%%%%%%%%%%%%%%%%%%%%%%%%%%%%%%%%%%%%%%%%%%%%%%%%%%%%
\section{Gauging}

%%%%%%%%%%%%%%%%%%%%%%%%%%%%%%%%%%%%%%%%%%%%%%%%%%%%%%%%%%%%%%%%
\frame{
\frametitle{Non-Abelian $p$-Form Hierarchy Structure} %Gauged Super-$p$-form Hierarchy
Have manifestly supersymmetric abelian $p$-form hierarchy.% and Chern-Simons-like invariant.
%
\mypause 
%
On $M=X\times Y$ there are mixed components of the frame $\sim dx^m e_m{}^i \partial/\partial y^i$ behaving like $G$ gauge fields for $G=\textrm{{\it Diff}} (Y)$
%
\mypause 
%
$\Rightarrow$ need to ``gauge'' the abelian tensor hierarchy.
%
\mypause 
%

Minimum structure required is a:\\ 
~\\
{\bf Non-abelian $p$-form hierarchy structure:}
$((\mathfrak g, f), (V_\bullet, t_\bullet, q),  h_\bullet)$ where
\begin{enumerate}
%
\mypause 
%
\item $(\mathfrak g, f)$ is a Lie algebra.
%
\mypause 
%
\item $(V_\bullet,t_\bullet, q)$ is a $\mathfrak g$-equivariant differential complex of representations $t_i : \mathfrak g\to GL(V_i)$.
%%
%\mypause 
%%
%\item For each $i$ a homomorphism of Lie algebras $t_i: \mathfrak g \to \mathrm{End} (V_i)$ defined by $t_i:a\mapsto t_a:= \rho_i(a)$.
%
\mypause 
%
\item For each $i$ a linear map $h_i: \mathfrak g \to \mathrm{Hom} (V_i,V_{i+1})$. 
\end{enumerate}
These data are subject to the following compatibility conditions: For all $a,b \in \mathfrak g$,
%\begin{enumerate}
%%
%\mypause 
%%
%\item Closure (intertwiner?) condition: $\forall a,b \in \mathfrak g$, %$[h_a , t_b] = h_{f(a,b)}$
%\begin{align*}
%[h_a , t_b] = h_{f(a,b)}
%\end{align*}
%%
%\mypause 
%%
%\item Nilpotence condition: $\forall a,b \in \mathfrak g$, %$h_a h_b + h_b h_a = 0$
%\begin{align*}
%h_a h_b + h_b h_a = 0
%\end{align*}
%%
%\mypause 
%%
%\item Cartan formula: $\forall a \in \mathfrak g$, %$q h_a + h_a q = t_a$
%\begin{align*}
%q h_a + h_a q = t_a
%\end{align*}
%\end{enumerate}
\begin{description}
%
\mypause 
%
\item[Closure:] $[h_a , t_b] = h_{f(a,b)}$
%\begin{align*}
%[h_a , t_b] = h_{f(a,b)}
%~~~
%\forall a,b \in \mathfrak g
%\end{align*}
%
\mypause 
%
\item[Nilpotence:] $h_a h_b + h_b h_a = 0$
%\begin{align*}
%h_a h_b + h_b h_a = 0
%\end{align*}
%
\mypause 
%
\item[Cartan formula:] $q h_a + h_a q = t_a$
%\begin{align*}
%q h_a + h_a q = t_a
%\end{align*}
\end{description}

%
\mypause 
%

Important example: Let $Y$ be a smooth manifold and set $(\mathfrak g, f)= (\mathfrak{X}(Y), [\cdot, \cdot])$ to be the algebra of diffeomorphisms and $(V_\bullet, q)=(\Omega_Y^\bullet, d_Y)$ to be the de Rham complex. Finally, for any vector field $a\in\mathfrak X(Y)$, set $t_a:=\mathscr L_a$ to be the Lie derivative along $a$ and set $h_a := \iota_a$ to be the contraction of forms by the vector field.

} %end frame


%%%%%%%%%%%%%%%%%%%%%%%%%%%%%%%%%%%%%%%%%%%%%%%%%%%%%%%%%%%%%%%%
\frame{
\frametitle{Construction of the Non-Abelian $p$-Form Hierarchy} %Gauged Super-$p$-form Hierarchy
We can now bolt this structure on to the superspace de Rham complex:
\begin{itemize}
%
\mypause 
%
\item In the abelian case, the $p$-form hierarchy structure can be used to make the double complex $(\mathscr O^\bullet_X(V_\bullet), d+q)$ of manifestly supersymmetric $p$-forms we already had with $V_i = \Omega^{p-i}_Y$ and $q=\partial$.
%
\mypause 
%
\item ``Minimal coupling'': Replace flat superspace $D\to \mathscr D$ with (super)curvature $\mathscr F \sim \mathscr D^2$.
%
\mypause 
%
\item Introduce derivation $\bm d := \mathscr D + q + h_\mathscr F$. 
%
\mypause 
%
\item By equivariance, $q^2=0$, closure relation, non-abelian Bianchi identity ($\mathscr D\mathscr F = 0$), and Cartan formula, $\bm d$ is a differential 
\begin{align*}
\bm d^2 = 0
\end{align*}
%
\mypause 
%
\item Repeat analysis of abelian super-$p$-form hierarchy: For $\mathbf A \in \bm \Omega^1$, $\bm \Lambda \in \bm \Omega^0$, {\it et cetera}
\begin{align*}
\delta \mathbf A = t_\lambda \mathbf A + \bm d \bm \Lambda
~~~,~~~
\mathbf F= \bm d \mathbf A
~~~\Rightarrow~~~
\delta \mathbf F = t_\lambda \mathbf F
~~~\mathrm{and}~~~
\bm d \mathbf F= 0.
\end{align*}
\end{itemize}
} %end frame

%%%%%%%%%%%%%%%%%%%%%%%%%%%%%%%%%%%%%%%%%%%%%%%%%%%%%%%%%%%%%%%%
\frame{
\frametitle{Non-Abelian Result}
Returning to the 11D supergravity example ({\it e.g.}\ $q\sim\partial$)
\begin{align*}
\begin{array}{ccc}
\begin{array}{ccl}
E &=&\hspace{8mm} - \,\partial \Phi\\
F &=& d_0(\Phi) - \partial V \\ 
W &=& d_1(V) - \partial \Sigma - \iota_\mathscr W \Phi \\ 
H &=& d_2(\Sigma) - \partial X  - \omega(\mathscr W, V) \\ 
G &=& d_3(X) \hspace{7mm} - \iota_\mathscr W \Sigma
\end{array}
&\Longrightarrow&
\begin{array}{ccl}
0&=& - \partial E\\
d_0(E)&=&- \partial F\\ 
d_1(F) &=& - \partial W - \iota_\mathscr W E \\ 
d_2(W) &=& - \partial H  - \omega(\mathscr W, F) \\ 
d_3(H) &=& - \partial G - \iota_\mathscr W W \\ 
d_4(G) &=& ~~~0
\end{array}
\end{array}
\end{align*}
The field strengths are all gauge covariant and satisfy the modified Bianchi identities.
%
\mypause 
%
The modified quadratic forms
\begin{align*}
\mathscr O^1 (\Omega_Y^7) \ni c_{[1,7]} &= \mathrm{Re}(E) \wedge F
\cr
\mathscr O^2 (\Omega_Y^6) \ni c_{[2,6]} &= - E  \wedge W -\tfrac i4 \bar {\mathscr D}^2 (F  \wedge {\mathscr D} F)
\cr
\mathscr O^3 (\Omega_Y^5) \ni c_{[3,5]} &= - \mathrm{Re}(E)  \wedge H -2 \omega (W, F) - i \iota_\mathscr W(F  \wedge {\mathscr D} F) + i \iota_{\bar {\mathscr W}} (F  \wedge \bar {\mathscr D} F)
\cr
\mathscr O^4 (\Omega_Y^4) \ni c_{[4,4]} &= E G + 2 W \wedge W -\tfrac i2 \bar {\mathscr D}^2 (FH)
\end{align*}
miraculously satisfy the modified descent relations.

} %end frame

%%%%%%%%%%%%%%%%%%%%%%%%%%%%%%%%%%%%%%%%%%%%%%%%%%%%%%%%%%%%%%%%
\frame{
\frametitle{Non-Abelian Chern-Simons Action}
To write an invariant action we need the analogue of a trace on the non-abelian $p$-form hierarchy structure. 
In the compactification case, this is $\int_Y$ with the conditions 
\begin{align*}
\int_Y \partial (\dots) = 0
~~~\mathrm{and}~~~
\int_Y \iota_\lambda (\dots) = 0 
~~~\Rightarrow~~~
\int_Y \mathscr L_\lambda (\dots) = 0 
%~~~~~ \lambda \in \mathfrak X(Y)
\end{align*}
The first is integration by parts. The second is non-existence of 8-forms on the 7-manifold $Y$. 
%(The third is diffeomorphism invariance.)

%
\mypause 
%


In a more general case, we need a 3-chain $\alpha$ on the non-abelian $p$-form hierarchy closed under the differentials $q$ and $h_\bullet$ (and therefore gauge-invariant). 
%
\mypause 
%
Explicitly, we need a chain
\begin{align*}
\alpha : \bigoplus_{r+s+t=4} \bm \Omega^r \otimes \bm \Omega^s \otimes \bm \Omega^t \to \mathbf R
\end{align*}
subject to 
\begin{align*}
q \alpha(\mathbf x, \mathbf y,\mathbf z) &:= 
	\alpha(q \mathbf a, \mathbf b,\mathbf c)
	+\alpha(\mathbf a, q\mathbf b,\mathbf c)
	+\alpha(\mathbf a, \mathbf b,q \mathbf c)
	=0
\cr	
h_a \alpha(\mathbf x, \mathbf y,\mathbf z)&:= 
	\alpha(h_a \mathbf x, \mathbf y,\mathbf z)
	+\alpha(\mathbf x, h_a \mathbf y,\mathbf z)
	+\alpha(\mathbf x, \mathbf y,h_a \mathbf z)
	=0
	~~~\forall a\in \mathfrak g
\end{align*}

} %end frame


%%%%%%%%%%%%%%%%%%%%%%%%%%%%%%%%%%%%%%%%%%%%%%%%%%%%%%%%%%%%%%%%
\frame{
\frametitle{Application: 11D Supergravity in $\mathbf R^{4|4} \times \mathbf R^7$}
Abstraction {\em necessary} to really describe higher-D theories in simple superspace:
%
\mypause 
%
\begin{itemize}
\item Cannot simply take $V_i = \Omega_Y^{7-i}$ because\\
$\Rightarrow$ decomposition of components in $Spin(7)$ representations\\
$\Rightarrow$ embedding in superfield representations gives ``fake'' components
%
\mypause 
%
%\item To single out 4D, $N=1$ superfield reps, need $\mathrm{Hol}(Y) = G_2$.
%\begin{align}
%\begin{array}{cccccc}
%g_{\bm m\bm n} &\to& g_{mn} & g_{mi} &g_{ij} \\
%\psi_{\bm m}^{\bm \alpha} &\to& \psi_m^{\alpha I} & \psi_j^{\alpha I} \\
%C_{\bm {mnp}} &\to& C_{mnp} & C_{mni} & C_{m ij} & C_{ijk}
%\end{array}
%\end{align}
%%
%\mypause 
%%
\item To single out 4D, $N=1$ superfield representations need $G_2 \subset Spin(7)$ decomposition:
\begin{align*}
\begin{array}{cccccccc}
&&\bm 1 & \bm 7& \bm 7 & \bm 1\oplus \bm {27} & \bm 7\oplus \bm {14} & \bm 1\oplus \bm {7}\oplus \bm {27}\\
g_{\bm m\bm n} &\to& g_{mn} & g_{mi} &&g_{ij} \\
\psi_{\bm m}^{\bm \alpha} &\to& \psi_m^\alpha&\psi_m^{\alpha i} & \eta_j^\alpha  &\chi_{(j}^{\alpha}{}_{i)}
	&\chi_{[j}^{\alpha}{}_{i]}  \\
C_{\bm {mnp}} &\to& C_{mnp} & C_{mni} &&& C_{m ij} & C_{ijk}
\end{array}
\end{align*}
%
\mypause 
%
\item 4D, $N=1$ component supermultiplets:
\begin{align*}
\begin{array}{ccl}
1 &:& (e_m^a, \psi_m^\alpha) \oplus ( g_{\bm 1} , C_{\bm 1}, C_{mnp} , \chi_{\bm 1}^\alpha) \\
7 &:& (C_{\bm 7m}, \psi_{\bm 7 m}^\alpha) \oplus (e_{\bm 7m}, \eta_{\bm 7}^\alpha) \oplus (C_{\bm 7}, C_{\bm 7 mn}, \chi_{\bm 7}^\alpha)\\
14 &:& (C_{\bm {14} m}, \chi_{\bm {14}}^\alpha) \\
27 &:& (C_{\bm {27}} + i g_{\bm {27}} , \chi_{\bm {27}}^\alpha) 
\end{array}
\end{align*}
\end{itemize}

} %end frame

%%%%%%%%%%%%%%%%%%%%%%%%%%%%%%%%%%%%%%%%%%%%%%%%%%%%%%%%%%%%%%%%
\frame{
\frametitle{Faithful Superspace Embedding}
\begin{itemize}
\item 1 superspin-$\tfrac32$ superconformal graviton $H^a \sim \dots + (\theta \sigma^m \bar \theta) e_m^a + \bar \theta^2 (\theta \sigma^m \psi_m)+\dots $
\item 7 superspin-$1$ gravitino superfields $\Psi^{\alpha i} \sim \dots + (\sigma^m\bar \theta)^\alpha C_{\bm 7 m} +(\theta \sigma^m \bar \theta) \psi_m^{\alpha i} + \dots$
\item 1 4-form ``radion'' field strength $G_{\bm 1} = (g_{\bm 1} + i C_{\bm 1}) + \theta \chi_{\bm 1} + \theta^2 \varepsilon^{mnpq} F_{mnpq} + \dots$
\item 7 3-form field strengths $H_{\bm 7} = C_{\bm 7} + \theta \chi_{\bm 7} + (\theta \sigma_m \bar \theta) \varepsilon^{mnpq} H_{npq}$
\item 21 2-form field strengths $W^\alpha_{\bm 7 \oplus \bm{14}} = \chi^\alpha_{\bm 7 \oplus \bm{14}} + (\theta \sigma^{mn})^\alpha (F_{{\bm 7 \oplus \bm{14}}}){}_{mn} +\dots$
\item 27 1-form field strengths $F_{\bm{27}} = g_{\bm{27}} + \theta \chi_{\bm {27}} + (\theta \sigma_m \bar \theta) \partial_m C_{\bm{27}}+ \dots$
\end{itemize}

%
\mypause 
%

Expanding in orders of the gravitational coupling we find:
\begin{itemize}
\item At quadratic order: Linearized ``old-minimal'' supergravity coupled to 7 gravitino multiplets $\bm \Psi_{\bm 7} = \Psi_{\bm 7} + W_{\bm 7} + \Sigma_{\bm 7}$ and the ``radion'' $G_{\bm 1}$. 
%
\mypause 
%
\item Cubic order contains a $p$-form hierarchy 
\begin{align*}
0\longrightarrow \mathscr O^0(\mathbf R^{27})
	\longrightarrow \mathscr O^1(\mathbf R^{7}\oplus \mathbf R^{14})
	\longrightarrow \mathscr O^2(\mathbf R^{7})
	\longrightarrow \mathscr O^3(\mathbf R)
	\longrightarrow 0
\end{align*}
%
\mypause 
%
\item Gauged by the non-abelian $V_{\bm 7}$ 
\end{itemize}


} %end frame

%%%%%%%%%%%%%%%%%%%%%%%%%%%%%%%%%%%%%%%%%%%%%%%%%%%%%%%%%%%%%%%%
%%%%%%%%%%%%%%%%%%%%%%%%%%%%%%%%%%%%%%%%%%%%%%%%%%%%%%%%%%%%%%%%
\section{Future work}
%%%%%%%%%%%%%%%%%%%%%%%%%%%%%%%%%%%%%%%%%%%%%%%%%%%%%%%%%%%%%%%%
\frame{
\frametitle{Open questions and current \& future work}
Loads of work to do... 
\begin{enumerate}
\item Relation to other ``non-abelian tensor hierarchies'', Courant Algebroids, Strongly homotopy Lie algebras, Lie algebra extensions, ... 
%%
%\mypause 
%%
\item Relation to on-shell $E_{7(7)}$ gauged supergravity and ``gravitational tensor hierarchies''
%%
%\mypause
%%
\item Relation to exceptional generalized geometry
%%
%\mypause
%%
\item Investigation of linearized action and resolution of ``the $\bm 7$-problem''
%%
%\mypause
%%
\item  Construction of non-linear superspace action, perhaps in gravitino expansion
%%
%\mypause
%%
\item  All orders in $\alpha'$ results for K\"ahler and super-potential
%%
%\mypause
%%
\item  ...
\end{enumerate}
%
\mypause
%
\begin{center}
{\bf \huge THANK YOU!}
\end{center}
}

%%%%%%%%%%%%%%%%%%%%%%%%%%%%%%%%%%%%%%%%%%%%%%%
{\tiny
%\small
%\footnotesize
%\linespread{1.1}\selectfont
\bibliography{/Users/wdlinch3/Dropbox/Rashoumon/LaTeX/BibTex/BibTex}
\bibliographystyle{unsrt}
} % end resize


\end{document}